
\documentclass{article}
\usepackage{colortbl}
\usepackage{makecell}
\usepackage{multirow}
\usepackage{supertabular}

\begin{document}

\newcounter{utterance}

\centering \large Interaction Transcript for game `mmlu\_pro', experiment `mmlu\_pro\_psychology', episode 257 with Qwen3.6{-}27B.
\vspace{24pt}

{ \footnotesize  \setcounter{utterance}{1}
\setlength{\tabcolsep}{0pt}
\begin{supertabular}{c@{$\;$}|p{.15\linewidth}@{}p{.15\linewidth}p{.15\linewidth}p{.15\linewidth}p{.15\linewidth}p{.15\linewidth}}
    \# & \multicolumn{2}{c}{Player} && \multicolumn{2}{c}{Game Master} \\
    \hline

    \theutterance \stepcounter{utterance}  
    & & & \multicolumn{4}{p{0.6\linewidth}}{
        \cellcolor[rgb]{0.9,0.9,0.9}{
            \makecell[{{p{\linewidth}}}]{
                \texttt{\tiny{[P1$\langle$GM]}}
                \texttt{The following are multiple choice questions (with answers) about psychology. Answer only with "The answer is (X)" where X is the correct letter choice.} \\
\\ 
\texttt{Question:} \\
\texttt{Pascale is interested in the processing strategies children use to learn new information. Pascale would best be classified as what type of psychologist?} \\
\texttt{Options:} \\
\texttt{A. social} \\
\texttt{B. school} \\
\texttt{C. sociocultural} \\
\texttt{D. forensic} \\
\texttt{E. behaviorist} \\
\texttt{F. health} \\
\texttt{G. clinical} \\
\texttt{H. cognitive} \\
\texttt{I. psychoanalytic} \\
\texttt{J. developmental} \\
\texttt{The answer is (H)} \\
\\ 
\texttt{Question:} \\
\texttt{According to Caplan's model of consultee{-}centered case consultation, the consultant is primarily interested in} \\
\texttt{Options:} \\
\texttt{A. identifying the causes and solutions of the client's presenting problems} \\
\texttt{B. establishing a hierarchy of authority to enable effective decision making} \\
\texttt{C. ensuring the consultee adheres strictly to a predetermined action plan} \\
\texttt{D. proposing multiple alternative solutions for the consultee to choose from} \\
\texttt{E. identifying the strengths and weaknesses of the consultee's current approach} \\
\texttt{F. presenting a single, well{-}defined and unambiguous course of action for the consultant to overcome skills deficits} \\
\texttt{G. developing a comprehensive treatment plan for the client} \\
\texttt{H. identifying and eliminating the causes of the consultee's difficulties in handling a problem} \\
\texttt{I. focusing on the consultant's personal growth and development} \\
\texttt{J. focusing on the relationship between the client and the consultee} \\
\texttt{The answer is (H)} \\
\\ 
\texttt{Question:} \\
\texttt{According to the Individuals with Disabilities Education Improvement Act, which of the following must an educational agency do before it changes the educational placement of a student with a disability?} \\
\texttt{Options:} \\
\texttt{A. Notify the parents in writing} \\
\texttt{B. Obtain the child's consent} \\
\texttt{C. Obtain a court order} \\
\texttt{D. Conduct a new evaluation of the child's disability} \\
\texttt{E. Discuss with the child's psychologist} \\
\texttt{F. Give the child a trial period in the new environment} \\
\texttt{G. Obtain parental consent} \\
\texttt{H. Notify the local education authority} \\
\texttt{I. Arrange a meeting with all teachers and administrators} \\
\texttt{J. Obtain school board approval} \\
\texttt{The answer is (A)} \\
\\ 
\texttt{Question:} \\
\texttt{Ani believes that her attitudes and behavior play a central role in what happens to her. Such a belief is likely to be associated with} \\
\texttt{Options:} \\
\texttt{A. low self{-}esteem.} \\
\texttt{B. a strong id.} \\
\texttt{C. a high level of anxiety.} \\
\texttt{D. a strong superego.} \\
\texttt{E. high levels of self{-}consciousness.} \\
\texttt{F. an external locus of control.} \\
\texttt{G. an inferiority complex.} \\
\texttt{H. a low level of self{-}awareness.} \\
\texttt{I. low self{-}efficacy.} \\
\texttt{J. an internal locus of control.} \\
\texttt{The answer is (J)} \\
\\ 
\texttt{Question:} \\
\texttt{In terms of Hofstede’s (1980) five cultural dimensions, the United States scores at the top of the scale on:} \\
\texttt{Options:} \\
\texttt{A. individualism and long{-}term orientation.} \\
\texttt{B. individualism and power distance.} \\
\texttt{C. uncertainty avoidance.} \\
\texttt{D. long{-}term orientation.} \\
\texttt{E. individualism.} \\
\texttt{F. individualism and masculinity.} \\
\texttt{G. long{-}term orientation and uncertainty avoidance.} \\
\texttt{H. power distance.} \\
\texttt{I. power distance and masculinity.} \\
\texttt{J. N/A} \\
\texttt{The answer is (E)} \\
\\ 
\texttt{Question:} \\
\texttt{Discuss the work cycle and its four human factors components.} \\
\texttt{Options:} \\
\texttt{A. Work cycle is a repetitive pattern of work, and its four human factors components are: anticipation processes, learning processes, coordination processes, and endurance processes.} \\
\texttt{B. Work cycle is the duration from one break to the next, and its four human factors components are: attention processes, memory consolidation, intuitive processes, and strength processes.} \\
\texttt{C. Work cycle is the process of planning and executing a series of tasks, and its four human factors components are: sensory perception, knowledge retrieval, creativity processes, and agility processes.} \\
\texttt{D. Work cycle is a series of tasks to be performed in a specific order, and its four human factors components are: sensory processes, memory processes, decision{-}making processes, and reflex processes.} \\
\texttt{E. Work cycle is the time it takes for a worker to return to the starting point of a task, and its four human factors components are: visual processing, auditory processing, strategic thinking, and precision processes.} \\
\texttt{F. Work cycle is the interval between the initiation and completion of a project, and its four human factors components are: situational awareness, mental modeling, risk assessment, and psychomotor functions.} \\
\texttt{G. Work cycle refers to the period of work before a break, and its four human factors components are: perceptual processes, recall processes, problem{-}solving processes, and emotional processes.} \\
\texttt{H. Work cycle is the sequence of operations performed by a machine or worker, and its four human factors components are: input processes, processing processes, output processes, and feedback processes.} \\
\texttt{I. Work cycle is the complete sequence of events necessary to complete a task, and its four human factors components are: perceptual processes, recall processes, problem solving (decision{-}making) processes, and muscular (motor) processes.} \\
\texttt{J. Work cycle is the sum of all tasks completed in a day, and its four human factors components are: cognitive processes, emotional processes, problem solving processes, and physical processes.} \\
            }
        }
    }
     \\ \\

    \theutterance \stepcounter{utterance}  
    & \multicolumn{4}{p{0.6\linewidth}}{
        \cellcolor[rgb]{0.561,0.737,0.561}{
            \makecell[{{p{\linewidth}}}]{
                \texttt{\tiny{[P1$\rangle$GM]}}
                \texttt{The answer is (I)} \\
            }
        }
    }
    & & \\ \\

    \theutterance \stepcounter{utterance}  
    & & & \multicolumn{2}{p{0.3\linewidth}}{
        \cellcolor[rgb]{0.9,0.9,0.9}{
            \makecell[{{p{\linewidth}}}]{
                \texttt{\tiny{[GM$|$GM]}}
                \texttt{I} \\
            }
        }
    }
    & & \\ \\

    \theutterance \stepcounter{utterance}  
    & & & \multicolumn{2}{p{0.3\linewidth}}{
        \cellcolor[rgb]{0.9,0.9,0.9}{
            \makecell[{{p{\linewidth}}}]{
                \texttt{\tiny{[GM$|$GM]}}
                \texttt{I} \\
            }
        }
    }
    & & \\ \\

    \theutterance \stepcounter{utterance}  
    & & & \multicolumn{2}{p{0.3\linewidth}}{
        \cellcolor[rgb]{0.9,0.9,0.9}{
            \makecell[{{p{\linewidth}}}]{
                \texttt{\tiny{[GM$|$GM]}}
                \texttt{game\_result = WIN} \\
            }
        }
    }
    & & \\ \\

\end{supertabular}
}

\end{document}
